\documentclass[12pt, a4paper]{article}
\usepackage[utf8]{inputenc}
\usepackage[T2A]{fontenc}
\usepackage{indentfirst, setspace}
\usepackage{tabularx, multirow}
\usepackage[normalem]{ulem}
\usepackage[style=russian]{csquotes}
\usepackage[english,russian]{babel}
\usepackage{hyperref}
\usepackage{ragged2e}
\usepackage{caption}
\usepackage{wrapfig}
\usepackage{amsmath}
\usepackage{tikz}
\makeatletter
\def\@biblabel#1{#1. }
\makeatother
\captionsetup{labelsep=endash}
\usepackage{listings}
\linespread{1.3}
\lstset{
  language=C++, 
  basicstyle=\ttfamily\small, 
  keywordstyle=\color{blue},
  breaklines=true,
  commentstyle=\color{green},
  stringstyle=\color{red}, 
  extendedchars=\true,
  showstringspaces=false,
}
\usepackage[left=20mm, top=15mm, right=15mm, bottom=20mm]{geometry}

\begin{document}
\begin{titlepage}
     \fontsize{12}{12}\selectfont

  {\centering

   \begin{bf}

    \begin{wrapfigure}{l}{19mm}
        
\includegraphics[width=22mm]{emblem.jpeg}
    \end{wrapfigure}


    \noindent Министерство науки и высшего образования Российской Федерации

    \noindent Федеральное государственное бюджетное образовательное учреждение высшего образования

    \noindent \enquote{Московский государственный технический унверситет

     \noindent имени Н.Э. Баумана

     \noindent (национальный исследовательский университет)}

    \noindent (МГТУ им. Н.Э. Баумана)

   \end{bf}
  }

  \vspace{0.4cm}

  {\setstretch{0.1}
   \noindent\rule{\textwidth}{1mm}
   \noindent\rule{\textwidth}{0.5mm}

  }

  \fontsize{14}{21}\selectfont

  \noindent\begin{tabularx}{\textwidth}{l >{\centering\arraybackslash}X}
   ФАКУЛЬТЕТ & \flqq Фундаментальные Науки\frqq \\ \cline{2-2}

   КАФЕДРА & ФН-12 \flqq Математическое моделирование\frqq \\ \cline{2-2}
  \end{tabularx}


  \vspace{1cm}


  \begin{center}
   \begin{bf}

    \fontsize{24}{36}\selectfont
    ОТЧЕТ

    \fontsize{20}{30}\selectfont
    ПО ЛАБОРАТОРНОЙ РАБОТЕ НА ТЕМУ:

    Сжатие файла методом Хаффмана

   \end{bf}
  \end{center}

  \fontsize{14}{21}\selectfont
  \vspace{5cm}


  \noindent\begin{tabularx}{\textwidth}{ X >{\centering}p{4cm} p{1cm} c }
   Студент: & & & Подобедов В. В. \\ \cline{2-2} \cline{4-4}
   & \fontsize{10}{15}\selectfont дата, подпись & & \fontsize{10}{15}\selectfont Ф.И.О. \\
   Преподаватель: & & & Волкова Л. Л.\\ \cline{2-2} \cline{4-4}
   & \fontsize{10}{15}\selectfont дата, подпись & & \fontsize{10}{15}\selectfont Ф.И.О.
  \end{tabularx}

  \vspace{\fill}

  \begin{center}
   \it{Москва}, 2023
  \end{center}

  \thispagestyle{empty}
\end{titlepage}

\renewcommand{\contentsname}{Содержание}
\setcounter{page}{2}
\tableofcontents

\newpage
\begin{center}
   \section*{Введение} 
\addcontentsline{toc}{section}{Введение}
\end{center}
\justifying
\textbf{Цель лабораторной работы}: реализовать алгоритмы сжатия по Хаффману.

Для достижения поставленной цели требуется решить следующие \textbf{задачи}.
\begin{enumerate}
\item Описать сжатие по Хаффману.
\item Реализовать сжатие файла.
\item Реализовать возможность считывать данные из файла.
\item Реализовать возможность записать ответ в бинарный файл.
\item Выполнить тестирование реализации разработанного алгоритма.
\end{enumerate}

Согласно варианту, требуется разработать сжатие файла по Хаффману для исходного кода программы на LISP.

\newpage
\begin{center}
   \section{Аналитическая часть} 
\end{center}

\textbf{Сжатие файла по алгоритму Хаффмана}~\cite{book} --- это метод кодирования данных, который использует переменные длины кодовых слов для представления символов в исходном тексте. Он основан на идее присвоения более коротких кодов более часто встречающимся символам и более длинных кодов менее часто встречающимся символам.

\textbf{LISP} --- это  функциональный язык программирования, который известен своим использованием списков как основной структуры данных и кода. Файлы LISP могут содержать определения функций, переменных, структур данных, а также вызовы функций, операции и другие выражения на языке LISP.

\newpage
\begin{center}
    \section{Конструкторская часть}
\end{center}

\subsection{Создание дерева Хаффмана}
\begin{enumerate}
    \item Подсчет частот символов.
    \begin{itemize}
        \item[---] Создать пустой словарь (map) ``freq'', где будут храниться частоты символов.
        \item[---] Пройтись по всем символам в тексте и увеличить частоту символа в словаре ``freq''. 
    \end{itemize}
    \item Создание приоритетной очереди.
        \begin{itemize}
        \item[---] Приоритетная очередь ``pq'' используется для хранения узлов дерева.
        \item[---] Для каждой буквы из словаря ``freq'' создать узел и поместить его в приоритетную очередь ``pq''. 
    \end{itemize}
    \item Построение дерева Хаффмана.
        \begin{itemize}
        \item[---] Выполнять нижеперечисленные действия, пока в приоритетной очереди не останется только один узел.
        \item[---] Извлечь два узла с минимальными частотами (две самые редкие буквы). Создать новый узел с суммой частот этих двух узлов и добавить обратно в очередь.
    \end{itemize}
    \item Получение корня дерева.
    \begin{itemize}
        \item[---] Корень дерева Хаффмана остается в приоритетной очереди после всех итераций. Извлечь его из очереди и присвоить его корню дерева Хаффмана.
    \end{itemize}
\end{enumerate}

\subsection{Кодирование элементов}
\begin{enumerate}
    \item Функция ``encode'' заполняет словарь ``huffmanCode'' кодами Хаффмана для каждой буквы.
    \item Параметры функции.
    \begin{itemize}
        \item[---] ``Node* root'' --- указатель на текущий узел дерева Хаффмана, начиная с корня.
        \item[---] ``string str'' --- строка, содержащая текущий код Хаффмана для данного узла.
        \item[---] ``map<char, string> \&huffmanCode''' --- словарь, в который мы записываем коды Хаффмана для символов. 
    \end{itemize}
    \item Проверка наличия узла.
    \begin{itemize}
        \item[---] Если текущий узел ``root'' равен ``nullptr'', то выйти из функции, так как достигнут конец ветви дерева.
    \end{itemize}
    \item Проверка листовых узлов.
    \begin{itemize}
        \item[---] Провеить, является ли текущий узел листовым, то есть у него нет потомков. Если это лист, записать текущий код ``str'' в словарь ``huffmanCode'' по ключу ``root->ch''. 
    \end{itemize}
    \item Рекурсивный обход дерева.
    \begin{itemize}
        \item[---] Рекурсивно вызвать ``encode'' для левого потомка текущего узла, добавить к текущему коду символа "$0$".
        \item[---] Рекурсивно вызвать ``encode'' для правого потомка текущего узла, добавить к текущему коду символа "$1$".
    \end{itemize}
\end{enumerate}

\subsection{Вычисление коэффициента сжатия}
\begin{enumerate}
    \item Подготовка переменных.
    \begin{itemize}
        \item[---] ``int bit\_1 = text.size() * 8'' рассчитать общее количество битов в исходном тексте. Умножение на 8 осуществляется, потому что каждый символ представляется 8 битами в стандартной кодировке ASCII.
        \item[---] ``int bit\_2 = 0'' эта переменная будет использоваться для подсчета общего количества битов в закодированном тексте.
    \end{itemize}
    \item Подсчет битов закодированного текста.
    \begin{itemize}
        \item[---] Цикл ``for'' проходит по каждой паре в словаре ``huffmanCode'', который содержит символы и их соответствующие коды Хаффмана.
        \item[---] Для каждой пары вычислить количество битов, занимаемых кодом Хаффмана для символа ``pair.second.size()'' и умножить на частоту символа в исходном тексте ``freq.at(pair.first)''. Это происходит потому, что сжатие зависит от того, сколько битов занимает каждый символ после кодирования и как часто он встречается в исходном тексте. Полученные значения присвоить переменной ``bit\_2''.
    \end{itemize}
    \item Вычисление коэффициента сжатия.
    \begin{itemize}
        \item[---] ``(bit\_1 + 0.0) / bit\_2'' Рассчитывает отношение количества битов исходного текста к количеству битов закодированного текста. Используется тип double, чтобы сохранить десятичную часть коэффициента.
    \end{itemize}
\end{enumerate}

\newpage
\begin{center}
    \section{Технологическая часть}
\end{center}

Лабораторная работа была реализована на языке программирования C++. Кроме стандартной библиотеки <iostream>, которая предоставляет возможность пользоваься средствами ввода-вывода для стандартной консоли, были подключена библиотека <queue>. Она представляет собой контейнер, который работает по принципу FIFO и обладает собственным набором функций. Была подключена библиотека <map> --- контейнер, представляющий отображение ключей на значения. Библиотека <vector> --- это  динамический массив, который обеспечивает быстрый доступ к элементам и динамическое изменение размера.Также была подключена библиотека <string> для пользования строками. Был также написан стек на ссылках. Для чтения и записи данных из/в файл была подключена библиотека <fstream>.

Далее приведен листинг реализации стека (листинг~\ref{lst:label1}).
\begin{lstlisting}[caption={Исходный код программы для сжатия по Хаффману},label={lst:label1}]
#include <iostream>
#include <fstream>
#include <queue>
#include <map>
#include <vector>
#include <string>

using namespace std;

struct Node {
	char ch;
	int freq;
	Node *left, *right;
};

struct Compare {
	bool operator()(Node* l, Node* r) {
		return (l->freq > r->freq);
	}
};

Node* getNode(char ch, int freq, Node* left, Node* right) {
	Node* node = new Node();
	node->ch = ch;
	node->freq = freq;
	node->left = left;
	node->right = right;
	return node;
}

void encode(Node* root, string str, map<char, string> &huffmanCode) {
	if (root == nullptr) return;
	if (!root->left && !root->right) {
		huffmanCode[root->ch] = str;
	}
	encode(root->left, str + "0", huffmanCode);
	encode(root->right, str + "1", huffmanCode);
}

void writeHuffmanCodesToFile(const map<char, string>& huffmanCode, const string& encodedText, const map<char, int>& freq) {
	ofstream outputFile("C:\\Users\\Kononova\\Desktop\\huffman_codes.bin", ios::binary);

	if (outputFile.is_open()) {
		for (const auto& pair : huffmanCode) {
			outputFile.put(pair.first);

			int frequency = freq.at(pair.first);
			outputFile.write(reinterpret_cast<const char*>(&frequency), sizeof(frequency));

			char codeLength = static_cast<char>(pair.second.length());
			outputFile.put(codeLength);
		}

		for (char ch : encodedText) {
			outputFile.put(ch);
		}

		outputFile.close();
		cout << "Данные успешно записаны в бинарный файл 'huffman_codes.bin'" << endl;
	}
	else {
		cout << "Невозможно открыть файл для записи" << endl;
	}
}

void winnings(string text, const map<char, string>& huffmanCode, const map<char, int>& freq) {
	int bit_1 = text.size() * 8;
	int bit_2 = 0;
	
	for (auto &pair : huffmanCode) {
		bit_2 += pair.second.size() *  freq.at(pair.first);
	}
	double win = (bit_1 + 0.0) / bit_2;
	cout << "Сжатие в " << win << " раз" << endl;
}

void buildHuffmanTree(string text) {
	map<char, int> freq;
	for (char ch : text) {
		freq[ch]++;
	}

	priority_queue<Node*, vector<Node*>, Compare> pq;	

	for (auto pair : freq) {
		pq.push(getNode(pair.first, pair.second, nullptr, nullptr));
	}

	while (pq.size() != 1) {
		Node* left = pq.top(); pq.pop();
		Node* right = pq.top(); pq.pop();

		int sumFreq = left->freq + right->freq;
		pq.push(getNode('\0', sumFreq, left, right));
	}

	Node* root = pq.top();

	map<char, string> huffmanCode;
	encode(root, "", huffmanCode);

	string encodedText;
	for (char ch : text) {
		encodedText += huffmanCode[ch];
	}

	cout << "Частоты букв: " << endl;
	for (auto pair : freq) {
		cout << pair.first << " : " << pair.second << endl;
	}
	cout << endl;

	cout << "Код букв: " << endl;
	for (auto pair : huffmanCode) {
		cout << pair.first << " : " << pair.second << endl;
	}
	cout << endl;

	cout << "Закодированный текст по Хаффману: " << encodedText << endl;

	writeHuffmanCodesToFile(huffmanCode, encodedText, freq);
	cout << endl;
	winnings(text, huffmanCode, freq);
}

int main() {
	setlocale(0, "");
	ifstream file("C:\\Users\\Kononova\\Desktop\\huff.lisp");
	string text = "";
	if (file.is_open()) {
		string text_help = "";
		while (getline(file, text_help)){
			text += text_help;
			text_help = "";
		}
		file.close();
	}
	else {
		cout << "Невозможно открыть файл";
		return 0;
	}
	cout << "Прочитанная фраза: " <<text << endl;
	buildHuffmanTree(text);
	
	return 0;
}
\end{lstlisting}


\newpage
\begin{center}
    \section{Исследовательская часть}
\end{center}

Далее на рисунках представлен результат работы программы. Размер входного файла с каждым тестом уменьшается на $20~\%$.

Результат работы программы, где из файла считывается программа на LISP, сжимает данные по алгоритму Хаффмана и записывает полученный ответ в бинарный файл.  Результат приведён на рис.~\ref{img:grap1}. 
\begin{figure}[h!]
    \center{\includegraphics[width=18cm]{test1_1.PNG}}
    \caption{Пример №1}
    \label{img:grap1}
\end{figure}
\newpage
\begin{figure}[h!]
    \center{\includegraphics[width=18cm]{test1_2.PNG}}
    \caption{Пример №1 (продолжение)}
\end{figure}

\newpage
Результат работы программы, где из файла считывается программа на LISP, сжимает данные по алгоритму Хаффмана и записывает полученный ответ в бинарный файл.  Результат приведён на рис.~\ref{img:grap2}. 
\begin{figure}[h!]
    \center{\includegraphics[width=18cm]{test2_1.PNG}}
    \caption{Пример №2}
    \label{img:grap2}
\end{figure}
\newpage
\begin{figure}[h!]
    \center{\includegraphics[width=18cm]{test2_2.PNG}}
    \caption{Пример №2 (продолжение)}
\end{figure}

\newpage
Результат работы программы, где из файла считывается программа на LISP, сжимает данные по алгоритму Хаффмана и записывает полученный ответ в бинарный файл.  Результат приведён на рис.~\ref{img:grap3}. 
\begin{figure}[h!]
    \center{\includegraphics[width=18cm]{test3_1.PNG}}
    \caption{Пример №3}
    \label{img:grap3}
\end{figure}
\newpage
\begin{figure}[h!]
    \center{\includegraphics[width=18cm]{test3_2.PNG}}
    \caption{Пример №3 (продолжение)}
\end{figure}

\newpage
Результат работы программы, где из файла считывается программа на LISP, сжимает данные по алгоритму Хаффмана и записывает полученный ответ в бинарный файл.  Результат приведён на рис.~\ref{img:grap4}. 
\begin{figure}[h!]
    \center{\includegraphics[width=18cm]{test4_1.PNG}}
    \caption{Пример №4}
    \label{img:grap4}
\end{figure}
\newpage
\begin{figure}[h!]
    \center{\includegraphics[width=18cm]{test4_2.PNG}}
    \caption{Пример №4 (продолжение)}
\end{figure}

\newpage
Результат работы программы, где из файла считывается программа на LISP, сжимает данные по алгоритму Хаффмана и записывает полученный ответ в бинарный файл.  Результат приведён на рис.~\ref{img:grap5}. 
\begin{figure}[h!]
    \center{\includegraphics[width=18cm]{test5_1.PNG}}
    \caption{Пример №5}
    \label{img:grap5}
\end{figure}
\begin{figure}[h!]
    \center{\includegraphics[width=18cm]{test5_2.PNG}}
    \caption{Пример №5 (продолжение)}
\end{figure}

\newpage
Результат работы программы, где из файла считывается программа на LISP, сжимает данные по алгоритму Хаффмана и записывает полученный ответ в бинарный файл.  Результат приведён на рис.~\ref{img:grap6}. 
\begin{figure}[h!]
    \center{\includegraphics[width=18cm]{test6.PNG}}
    \caption{Пример №6}
    \label{img:grap6}
\end{figure}

\newpage
Результат работы программы, где из файла считывается программа на LISP, сжимает данные по алгоритму Хаффмана и записывает полученный ответ в бинарный файл.  Результат приведён на рис.~\ref{img:grap7}. 
\begin{figure}[h!]
    \center{\includegraphics[width=18cm]{test7.PNG}}
    \caption{Пример №7}
    \label{img:grap7}
\end{figure}

\newpage
Результат работы программы, где из файла считывается программа на LISP, сжимает данные по алгоритму Хаффмана и записывает полученный ответ в бинарный файл.  Результат приведён на рис.~\ref{img:grap8}. 
\begin{figure}[h!]
    \center{\includegraphics[width=18cm]{test8.PNG}}
    \caption{Пример №8}
    \label{img:grap8}
\end{figure}

\newpage
Результат работы программы, где из файла считывается программа на LISP, сжимает данные по алгоритму Хаффмана и записывает полученный ответ в бинарный файл.  Результат приведён на рис.~\ref{img:grap9}. 
\begin{figure}[h!]
    \center{\includegraphics[width=18cm]{test9.PNG}}
    \caption{Пример №9}
    \label{img:grap9}
\end{figure}

\newpage
Результат работы программы, где из файла считывается программа на LISP, сжимает данные по алгоритму Хаффмана и записывает полученный ответ в бинарный файл.  Результат приведён на рис.~\ref{img:grap10}. 
\begin{figure}[h!]
    \center{\includegraphics[width=18cm]{test10.PNG}}
    \caption{Пример №10}
    \label{img:grap10}
\end{figure}

Анализируя выполненные тесты, можно отметить, что при увеличении объема данных алгоритм Хаффмана становится более эффективным в сжатии по сравнению с простым методом сжатия. Это происходит благодаря созданию оптимального кода для представления данных: наиболее часто встречающиеся символы кодируются более короткими кодами, что позволяет значительно сократить общий размер данных.

Однако, когда данные становятся более случайными или каждый символ встречается с примерно одинаковой частотой, эффективность сжатия по Хаффману может уменьшаться. Это происходит из-за того, что оптимальное кодирование для равномерно распределенных данных не обеспечивает большого преимущества в сжатии.

Таким образом, при увеличении объема данных сжатие по Хаффману обычно стремится к более эффективному сжатию, но его эффективность может варьироваться в зависимости от характера данных и распределения символов в них.

\newpage
\begin{center}
    \section*{Заключение}
    \addcontentsline{toc}{section}{Заключение}
\end{center}

Цель лабораторной работы достигнута: реализован алгоритм сжатия по Хаффману.

Для достижения поставленной цели требуется решить следующие \textbf{задачи}.
\begin{enumerate}
\item Описано сжатие по Хаффману.
\item Реализовано сжатие файла.
\item Реализована возможность считывать данные из файла.
\item Реализована возможность записать ответ в бинарный файл.
\item Выполнено тестирование реализации разработанного алгоритма.
\end{enumerate}

\newpage
\begin{center}
\begin{thebibliography}{}
\addcontentsline{toc}{section}{Список используемых источников}
\bibitem{book} Котиева, Х. М. Сжатие данных без потерь. Использование алгоритма Хаффмана / Х. М. Котиева. — Текст : непосредственный // Молодой ученый. — 2020. — № 35 (325 с.).
\end{thebibliography}
\end{center}

\end{document}